% Generated by analysis/make_tables.py -- do not edit by hand.

% needs the packages array, booktabs and longtable, and \ProvenanceCaption
\begin{longtable}{@{}>{\raggedright\arraybackslash}p{0.04\linewidth}>{\raggedright\arraybackslash}p{0.19\linewidth}>{\raggedright\arraybackslash}p{0.07\linewidth}>{\raggedright\arraybackslash}p{0.21\linewidth}>{\raggedright\arraybackslash}p{0.08\linewidth}>{\raggedright\arraybackslash}p{0.07\linewidth}>{\raggedright\arraybackslash}p{0.16\linewidth}@{}}
\caption{\ProvenanceCaption}\\
\toprule
ID & Quantity & Kind & Source & Location & Units & Rule of use \\
\midrule
\endfirsthead
\toprule
ID & Quantity & Kind & Source & Location & Units & Rule of use \\
\midrule
\endhead
\bottomrule
\endlastfoot
P01 & $C_{11}, C_{12}, C_{44}$ of Al, Cu, Fe, W (2NN MEAM) & published & Lazar, Agiasofitou and Boehlke (2022)~\cite{Lazar2022}; computed with the 2NN MEAM potential of Lee et al., as stated there (OpenKIM MD\_111291751625\_001), by the method of Admal et al. (2017)~\cite{Admal2017} & Table 1, p. 117 & eV/\AA$^3$ & Input. Al and Cu MEAM for all results, Fe and W for stiffness and Voigt only (D not positive definite); for Al and Cu each entry from the finer of P01 and P06 (analysis/inputs.py) \\
P02 & $D_{1,1},\dots,D_{16,17}$ of Al, Cu, Fe, W (2NN MEAM) & published & Lazar, Agiasofitou and Boehlke (2022)~\cite{Lazar2022}; same potential and method as P01 & Table 1, p. 117 & eV/\AA & Input; for Al and Cu each entry from the finer of P02 and P06 (analysis/inputs.py); $a_i$ are recomputed from these values (never copied from Table 2) \\
P03 & $a_1,\dots,a_{11}$ of Al, Cu, Fe, W & published & Lazar, Agiasofitou and Boehlke (2022)~\cite{Lazar2022} & Table 2, p. 117 & eV/\AA & Cross-check only; not an input \\
P04 & $C_{11}, C_{12}, C_{44}$ of Cu (EAM) & published & Po, Admal and Lazar (2019)~\cite{Po2019}; EAM potential of Mendelev et al. (2008) according to their text; method of Admal et al. & Table 1, p. 11 & eV/\AA$^3$ & Input (third cubic data set; Po et al. note possible artefacts of EAM gradient constants) \\
P05 & $D_{1,1},\dots,D_{16,17}$ of Cu (EAM) & published & Po, Admal and Lazar (2019)~\cite{Po2019}; as P04 & Table 1, p. 11 & eV/\AA & Input \\
P06 & C and D of Cu and Al (2NN MEAM), second print of P01-P02 & published & Po, Admal and Lazar (2019)~\cite{Po2019} & Table 1, p. 11 & eV/\AA$^3$; eV/\AA & Input where its digits are finer than those of P01-P02: D2,5, D3,5, D16,16, D16,17 of Cu and D2,5 of Al (analysis/inputs.py); elsewhere the same digits or fewer \\
P07 & C and $a_1,\dots,a_{11}$ of Al, Cu, Fe, W (repeat of Lazar et al. 2022) & published & Lazar and Po (2024)~\cite{Lazar2024} & Table 1, p. 23 & eV/\AA$^3$; eV/\AA & Cross-check only; not an independent data set \\
P08 & characteristic length $\ell$ of Al, Cu, Fe, W & published & Lazar and Po (2024)~\cite{Lazar2024} & Table 2, p. 24 & \AA & Comparison value: reproduced for Al and Cu; for Fe and W their printed $\ell$ differs in the last digits from their $\ell$/a times their a, which is reproduced \\
P09 & lattice parameter a of Al, Cu, Fe, W & published & Lazar and Po (2024)~\cite{Lazar2024} (their source not stated) & Table 2, p. 24 & \AA & Context only (grain size relative to the lattice parameter; $\ell$/a; nearest-neighbour distance b = a $\sqrt3$/2 of Fe and W) \\
P10 & $\ell$/a of Al, Cu, Fe, W & published & Lazar and Po (2024)~\cite{Lazar2024} & Table 2, p. 24 & none & Cross-check only \\
P11 & Voigt-averaged isotropic constants $\bar a_1,\dots,\bar a_5$ of Al and W & published & Lazar, Agiasofitou and Boehlke (2022)~\cite{Lazar2022}; the same values are in Table 1 of Eremeyev and Lazar (2022)~\cite{Eremeyev2022} & Table 6, p. 132 & eV/\AA & Comparison value: our Voigt estimate reproduces it to $7\times10^{-6}$ \\
P12 & $C_{11}, C_{12}, C_{44}$ of Solyaev's Fe grains & published & Solyaev (2023)~\cite{Solyaev2023}, citing Lubarda (his ref. [8]) & Section 5, p. 4 & GPa & Supplementary Fe case only; not a validation of the implementation \\
P13 & bulk modulus $k$ = ($C_{11}+2C_{12}$)/3 and classical Poisson ratio of the Fe polycrystal & published & Solyaev (2023)~\cite{Solyaev2023} & Section 5, p. 4 & GPa; none & Comparison values for the supplementary Fe case \\
P14 & $\ell$/$\ell_0$ and $\nu$ versus $2R/\ell_0$; minima of the error function & derived & Solyaev (2023)~\cite{Solyaev2023}, digitised in this work & Figs. 1-2, p. 5 & none & Supplementary Fe case only \\
P15 & evaluation point ($\nu$ = 0.348, $\ell$ = 1.9 $\ell_0$, R = 5 $\ell_0$) and grain sizes 2R = 10, 50 $\ell_0$ & assumed & this work (point chosen near Solyaev's printed solution; sizes of his Fig. 1 minima) & this work & none ($\ell_0=1$) & Supplementary Fe case \\
P16 & fixed grid of medium length and grain radius for the weight of the sixth condition & assumed & this work & this work & \AA & Diagnostic only \\
P17 & contact stiffnesses $k_n$, $k_s$, fabric parameter $a_{20}$ of the transversely isotropic packing and $\Delta/\sigma$ sweep of the granular hexagonal tensors & assumed & this work, within the range of figure 3 of Pirmoradi, Suiker and Poorsolhjouy (2024)~\cite{Pirmoradi2024} ($k_s$ = $k_n$/2, 0 $\le$ $a_{20}$ $\le$ 0.2); tensors by their eqs. (52), (70) and (71) & Pirmoradi et al. 2024, Fig. 3, p. 675; eqs. (70)-(71), pp. 668-669 & none ($\rho_c l^4k_n=1$) & Model inputs of the granular example; the extra modulus follows from $a_{66}$ and is swept algebraically beyond the reachable range \\
P18 & grain radii and the starting grid of the minimisation & assumed & this work & this work & \AA{} ($R$, $\ell$); none ($\nu$) & Numerical settings; every local minimum is polished \\
P19 & $C_{11}, C_{12}, C_{44}$ of a strongly anisotropic Fe-like grain (Zener ratio 2.46) and the grain sizes $2R/\ell_0$ = 10, 20, 50, 100 of its runs and of the mixed readings & assumed & this work & this work & GPa ($C$); none ($2R/\ell_0$) & Illustrative only: shows that a strong anisotropy still gives l close to $\ell_0$; not a material \\
\end{longtable}
